\documentclass{article}
% if you need to pass options to natbib, use, e.g.:
%     \PassOptionsToPackage{numbers, compress}{natbib}
% before loading neurips_2023
\usepackage[final]{neurips}
% to avoid loading the natbib package, add option nonatbib:
%    \usepackage[nonatbib]{neurips}


\usepackage[utf8]{inputenc} % allow utf-8 input
\usepackage[T1]{fontenc}    % use 8-bit T1 fonts
\usepackage{hyperref}       % hyperlinks
\usepackage{url}            % simple URL typesetting
\usepackage{booktabs}       % professional-quality tables
\usepackage{amsfonts}       % blackboard math symbols
\usepackage{nicefrac}       % compact symbols for 1/2, etc.
\usepackage{microtype}      % microtypography
\usepackage{xcolor}         % colors
\usepackage{graphicx}       % for including figures


\title{Project Title}


% The \author macro works with any number of authors. There are two commands
% used to separate the names and addresses of multiple authors: \And and \AND.
%
% Using \And between authors leaves it to LaTeX to determine where to break the
% lines. Using \AND forces a line break at that point. So, if LaTeX puts 3 of 4
% authors names on the first line, and the last on the second line, try using
% \AND instead of \And before the third author name.


\author{%
  Name\\
  Department of ...\\
  University of Toronto\\
  \texttt{email} \\
  % more authors
  \And
  Name\\
  Department of ...\\
  University of Toronto\\
  \texttt{email} \\
  % \AND
  % Coauthor \\
  % Affiliation \\
  % Address \\
  % \texttt{email} \\
  % \And
  % Coauthor \\
  % Affiliation \\
  % Address \\
  % \texttt{email} \\
  % \And
  % Coauthor \\
  % Affiliation \\
  % Address \\
  % \texttt{email} \\
}

\begin{document}


\maketitle


\begin{abstract}
An abstract of the report.
\end{abstract}

\section*{Attestation of Teamwork}
Here, you explain how each member in the group has contributed to the project. 



\section{Introduction}
Here, you explain briefly the target problem and motivation to solve it via Deep Learning.

\section{Preliminaries and Problem Formulation}
Here, you formulate the problem. Specify the ultimate goal. If there are some new concepts that you had to learn to understand this project, you could briefly explain them. 
	
\section{Design}

% ── SVM SUBSECTION (Jake) ────────────────────────────────────────────────
\subsection{SVM-Based Intrusion Detection}

We design an SVM-based classifier for IEC~104 SCADA network intrusion detection that distinguishes normal traffic from six attack types (flood, fuzzy, MITM, IEC-104 starvation, NTP DDoS, and port scan). The system follows a three-phase improvement pipeline:

\textbf{Phase~1 (Baseline)} uses a linear-kernel SVM with default parameters on raw, unscaled features. This establishes a performance floor by deliberately omitting scaling, non-linear kernels, and class balancing.

\textbf{Phase~2 (Improved)} addresses each baseline weakness: StandardScaler normalises features to zero-mean unit-variance, ANOVA F-test selects the 30 most discriminative features, an RBF kernel captures non-linear decision boundaries, and \texttt{class\_weight=`balanced'} up-weights minority classes inversely proportional to their frequency.

\textbf{Phase~3 (Optimised)} performs exhaustive hyperparameter tuning via GridSearchCV over $C \in \{0.1, 1, 10, 100\}$, $\gamma \in \{\text{scale, auto}, 0.01, 0.1\}$, and kernel $\in \{\text{rbf, poly}\}$, evaluated with 5-fold stratified cross-validation scored by weighted-F1.

% ── PLACEHOLDER FOR TEAMMATE'S MODEL ────────────────────────────────────
% \subsection{[Teammate's Model] Design}
% [Teammate adds their model design here]


\section{Methodology}

% ── SVM SUBSECTION (Jake) ────────────────────────────────────────────────
\subsection{SVM Methodology}

\textbf{Data Preprocessing.}
The dataset comprises 7 CSV files from CICFlowMeter containing 84 columns per flow. We drop four metadata columns (Flow~ID, Src/Dst~IP, Timestamp) that encode host identity rather than traffic patterns, replace infinite values and NaNs with column medians, and encode the target label via \texttt{LabelEncoder}. A 75/25 stratified train-test split preserves class proportions.

\textbf{Feature Selection.}
ANOVA F-test ranks features by their between-class to within-class variance ratio. We retain the top 30 features; the highest-scoring features (SYN Flag Count, Fwd PSH Flags, FIN Flag Count) align with known IEC~104 attack signatures where attackers produce anomalous TCP flag patterns.

\textbf{Multi-Seed Evaluation.}
To reduce sensitivity to a single train-test split, we run the full three-phase pipeline with five different random seeds ($\{42, 123, 256, 789, 1024\}$) and report the averaged metrics with standard deviations. Each seed controls the stratified split, the SVM's internal randomisation, and the cross-validation fold assignment.

\textbf{Implementation.}
The system is implemented in Python using scikit-learn (SVC, GridSearchCV, StandardScaler, SelectKBest), pandas for data handling, and matplotlib/seaborn for visualisation. All experiments use \texttt{n\_jobs=-1} for parallelised grid search.

% ── PLACEHOLDER FOR TEAMMATE'S MODEL ────────────────────────────────────
% \subsection{[Teammate's Model] Methodology}
% [Teammate adds their methodology here]


\section{Numerical Experiments}

\subsection{SVM Results}

Table~\ref{tab:svm-results} presents the averaged performance metrics across five random seeds for each improvement phase.

\begin{table}[h]
  \caption{SVM performance across improvement phases (mean $\pm$ std over 5 seeds)}
  \label{tab:svm-results}
  \centering
  \small
  \begin{tabular}{lccc}
    \toprule
    Metric & Phase 1 (Baseline) & Phase 2 (Improved) & Phase 3 (Optimised) \\
    \midrule
    Accuracy        & $0.964 \pm 0.002$ & $0.999 \pm 0.000$ & $\mathbf{1.000 \pm 0.000}$ \\
    F1 (weighted)   & $0.957 \pm 0.003$ & $0.999 \pm 0.000$ & $\mathbf{1.000 \pm 0.000}$ \\
    F1 (macro)      & $0.449 \pm 0.022$ & $0.894 \pm 0.015$ & $\mathbf{0.933 \pm 0.010}$ \\
    Precision (wt)  & $0.959 \pm 0.003$ & $1.000 \pm 0.000$ & $\mathbf{1.000 \pm 0.000}$ \\
    Recall (wt)     & $0.964 \pm 0.002$ & $0.999 \pm 0.000$ & $\mathbf{1.000 \pm 0.000}$ \\
    \bottomrule
  \end{tabular}
\end{table}

The baseline achieves 96.4\% accuracy but only 44.9\% macro-F1, indicating near-complete failure on minority classes. Phase~2 delivers the largest improvement: adding scaling, feature selection, RBF kernel, and class balancing lifts macro-F1 from 0.449 to 0.894. Phase~3's GridSearchCV further refines performance to 93.3\% macro-F1 and perfect weighted accuracy, selecting a polynomial kernel with $C\!=\!100$ and $\gamma\!=\!0.1$ as the optimal configuration.

Figure~\ref{fig:svm-comparison} shows the performance trajectory across phases and the F1 (macro) improvement curve. Figure~\ref{fig:svm-confusion} presents normalised confusion matrices for all three phases.

\begin{figure}[h]
  \centering
  \includegraphics[width=\textwidth]{fig05_performance_comparison.png}
  \caption{SVM performance metrics across phases (left) and F1 macro improvement trajectory (right).}
  \label{fig:svm-comparison}
\end{figure}

\begin{figure}[h]
  \centering
  \includegraphics[width=\textwidth]{fig06_confusion_matrices.png}
  \caption{Normalised confusion matrices (\% of true class) for each SVM phase. Diagonal values indicate correct classification rate.}
  \label{fig:svm-confusion}
\end{figure}

% ── PLACEHOLDER FOR TEAMMATE'S RESULTS ──────────────────────────────────
% \subsection{[Teammate's Model] Results}
% [Teammate adds their results, figures, and tables here]


\section{Discussion}

\subsection{SVM Analysis}

\textbf{Impact of preprocessing.}
The dominant improvement comes from Phase~2, not Phase~3. Feature scaling alone is critical for SVMs because the RBF kernel's distance computation is dominated by high-range features (e.g., Flow Bytes/s $\sim$10\textsuperscript{7}) without normalisation. Class balancing via inverse-frequency weighting was equally important --- without it, the SVM effectively ignores classes with fewer than 100 samples.

\textbf{Remaining challenges.}
MITM attack detection remains the weakest point (F1 $\approx$ 0.57--0.67 across seeds), attributable to its extremely small sample size (26 flows). Flood attack detection also shows higher variance across seeds. These results suggest that data augmentation or specialised sampling strategies could further improve minority-class performance.

\textbf{Multi-seed robustness.}
The low standard deviations across five seeds ($\leq$0.022 for all metrics) confirm that the results are not artefacts of a single fortunate split. The optimised model consistently achieves near-perfect weighted metrics regardless of the random initialisation.

% ── PLACEHOLDER FOR TEAMMATE'S DISCUSSION ───────────────────────────────
% \subsection{[Teammate's Model] Analysis}
% [Teammate adds their discussion here]

	

\section{Conclusions}
You should shortly conclude what you learned during project. Also, if you have any idea for further improvement of the results or extension, you may mention it here.
	
	


\section*{References}
Include all references here. It's important to have your references cited.
\medskip
{
\small


[1] Alexander, J.A.\ \& Mozer, M.C.\ (1995) Template-based algorithms for
connectionist rule extraction. In G.\ Tesauro, D.S.\ Touretzky and T.K.\ Leen
(eds.), {\it Advances in Neural Information Processing Systems 7},
pp.\ 609--616. Cambridge, MA: MIT Press.


[2] Bower, J.M.\ \& Beeman, D.\ (1995) {\it The Book of GENESIS: Exploring
  Realistic Neural Models with the GEneral NEural SImulation System.}  New York:
TELOS/Springer--Verlag.


[3] Hasselmo, M.E., Schnell, E.\ \& Barkai, E.\ (1995) Dynamics of learning and
recall at excitatory recurrent synapses and cholinergic modulation in rat
hippocampal region CA3. {\it Journal of Neuroscience} {\bf 15}(7):5249-5262.
}

\section*{Appendix}

\subsection*{A. SVM Supplementary Figures}

\begin{figure}[h]
  \centering
  \includegraphics[width=0.85\textwidth]{fig01_class_distribution.png}
  \caption{Class distribution in the IEC 104 dataset, highlighting severe imbalance (dosattack: 304,627 vs mitmattack: 26).}
  \label{fig:app-class-dist}
\end{figure}

\begin{figure}[h]
  \centering
  \includegraphics[width=0.85\textwidth]{fig04_feature_importance.png}
  \caption{Top 30 features ranked by ANOVA F-test score. SYN Flag Count, Fwd PSH Flags, and FIN Flag Count are the most discriminative.}
  \label{fig:app-features}
\end{figure}

\begin{figure}[h]
  \centering
  \includegraphics[width=\textwidth]{fig07_per_class_f1.png}
  \caption{Per-class F1 score across improvement phases. Largest gains in flood (0.42$\to$0.91), fuzzy (0.05$\to$1.00), and NTP DDoS (0.01$\to$1.00).}
  \label{fig:app-perclass}
\end{figure}

\begin{figure}[h]
  \centering
  \includegraphics[width=0.7\textwidth]{fig09_gridsearch_heatmap.png}
  \caption{GridSearchCV heatmap showing weighted-F1 for the polynomial kernel across C and gamma values.}
  \label{fig:app-gridsearch}
\end{figure}

\begin{figure}[h]
  \centering
  \includegraphics[width=0.7\textwidth]{fig10_learning_curve.png}
  \caption{Learning curve for the optimised SVM. Training and cross-validation scores converge near 1.0, indicating no overfitting.}
  \label{fig:app-learning}
\end{figure}


\newpage
\section*{AI Usage}
Please explain in one paragraph how you have used AI to accomplish this project. If you have not used AI at all, please indicate that you have not used it. 

\section*{Contest for Information Sharing}
As part of this course, we may share selected project materials (e.g., reports, presentation slides, and presentation recordings) on the course webpage as learning resources for future students. Additionally, we may use anonymized project information for internal statistical analysis of course outcomes. Please indicate your preferences below. \textit{Your choices will not affect your grade in any way.}

\subsection*{Consent for Sharing Project Materials}
\textit{Please keep the item you wish and remove the other one.}

\begin{itemize}
	\item All group members consent to allow the project materials (report, slides, and presentation recording) to be shared on the course page for future students.
	\item The group members do \textbf{not} consent to allow the project materials (report, slides, and presentation recording) to be shared on the course page for future students.
\end{itemize}

\subsection*{Consent for Use of Project Information in Statistical Analysis}
\textit{Please keep the item you wish and remove the other one.}
\begin{itemize}
	\item All group members consent to allow anonymized information about the project (e.g., topic, methods, outcomes, grades) to be used by the instructor for statistical analysis and course improvement.
	\item The group members do \textbf{not} consent to allow anonymized information about the project (e.g., topic, methods, outcomes, grades) to be used by the instructor for statistical analysis and course improvement.
\end{itemize}

\subsection*{Optional Comments}
\textit{Please list any specific conditions or comments here.}

\subsection*{Group Identification}

\begin{itemize}
	\item Group number:
	\item Names of group members: 
	\item Signature of of group members: 
	\item Date: 
\end{itemize}

\end{document}
